\documentclass[../main.tex]{subfiles}
\begin{document}
%==========================================TIÊU ĐỀ TẬP========================================
\begin{table}[h]
  \begin{adjustwidth}{-1cm}{}
    \begin{tabular}{>{\centering\arraybackslash}p{6cm}|>{\raggedright\arraybackslash}p{12.5cm}}
      \multirow{5}{6cm}
      % Cột bên trái: ÉPISODE và số tập
      {\\[-40pt]
        \flaregothic\fontsize{20pt}{20pt}\selectfont \centering
        \color{eptype}ÉPISODE\color{black}\\
        \burbank\fontsize{72pt}{72pt}\selectfont
        \color{epnum}57\color{black}\\[6pt]
        \cafeteria\fontsize{20pt}{20pt}\selectfont\color{epnum}(5-1)\color{black}
        \\[18pt]
        \color{black}
      }
       &
      % Dòng thứ nhất: tên tập tiếng Nhật
      {\fotskip\fontsize{18pt}{18pt}\selectfont \ruby{東}{とう}\ruby{大}{だい}\footnotemark を\ruby{目}{め}\ruby{指}{ざ}して
      } \\[6pt]
      % Dòng thứ hai: tên tập tiếng Anh
       & {\cafeteria\fontsize{18pt}{18pt}\selectfont Akidemic Pursuits} \\[4pt]
      % Dòng thứ ba: tên tập tiếng Việt
       & {\vnmsans\fontsize{18pt}{18pt}\selectfont Mọt sách Aki - Nhắm vào đại học Tokyo!} \\[8pt]
      % Dòng thứ tư: Nhân vật xuất hiện
       & {\caviar{\fontsize{14pt}{14pt}\selectfont Nhân vật: \emp{kuon1} \emp{suisei} \emp{aki} \emp{matsuri} \emp{marine}}}
      \\[8pt]
      % Dòng cuối cùng: Ngày phát hành
       & {\continuum\fontsize{16pt}{16pt}\selectfont Ngày phát hành: 07/06/2020} \\[18pt]
    \end{tabular}
  \end{adjustwidth}
\end{table}
\footnotetext{Viết tắt của 東京大学 (とうきょうだいがく), hay Toudai, tức đại học Tokyo.}
%=========================================TRUYỆN CHÍNH========================================
\color{black}

\txtbx[2]{{\color{vol} \uline{Tóm tắt:}} Aki bỗng dưng biến thành ``mọt sách'' cuồng học, miệng lẩm bẩm ``ba nhân ba bằng Chòm Nam Thập Tự''. Marine, Suisei và Matsuri dùng đủ mọi chiêu: bánh ngọt, bóng đá, mời đi chơi - đều thất bại. Suisei và Matsuri quyết định học giỏi hơn để cạnh tranh, nhưng kết quả là cả hai hóa thành\dots\ báo đốm. Aki thuyết giảng về loài báo với những ``sự thật'' nổ não: chúng sống ở Akihabara, ăn spaghetti bay, và đặt câu hỏi về chiến tranh. Kết cục, Kuon phải gọi bệnh viện tâm thần cho Aki, còn hai con báo đốm thì bị lôi về nhà thuần hóa.}

Tháng Sáu gõ cửa, mang theo cái nắng vàng óng ả của những ngày hè đầu tiên. Trời xanh trong vắt, không một gợn mây, chỉ có những tia nắng nhảy múa trên từng tán cây ngoài cửa sổ. Tiếng ve râm ran như một dàn đồng ca bất tận, giục giã, nhắc nhở rằng mùa thi đang cận kề, nhưng cũng báo hiệu một mùa hè sôi động sắp bắt đầu.

Tuy nhiên, bên trong văn phòng \hll, bầu không khí lại chẳng hề rộn ràng. Nó ngột ngạt, im lìm, và kỳ quặc đến lạ thường.

Ở chiếc bàn gần cửa sổ, Aki đang ngồi học bài, một điều mà cô hầu như không bao giờ làm. Trông cô lúc này chẳng khác gì một sĩ tử thực thụ: lưng thẳng, mắt chăm chú, tay cầm bút lia lịa. Trước mặt cô là một chồng sách cao ngất, đủ để làm một chiếc thang leo lên tầng hai. Bút chì, vở bài tập, giấy nháp rải rác khắp bàn như một chiến trường sau trận đánh. Điều đó cũng phải thôi; đây là thời điểm mà toàn đất nước mặt trời mọc đang chuẩn bị cho kỳ thi đại học sắp tới, nên có lẽ cô nàng đang muốn hưởng ứng với tất cả mọi người.

Thỉnh thoảng, cô lẩm bẩm điều gì đó; có vẻ là công thức toán học, nhưng cũng có thể là tên của các vì sao, hoặc một thứ ngôn ngữ nào đó không ai hiểu.

\img{5-1a}

Hành động kỳ quặc này không thể thoát khỏi sự chú ý của những người xung quanh. Matsuri, Marine, Suisei và Kuon tụ lại ở cửa phòng, đứng thành một cụm, mắt mở to, nhìn chằm chằm vào nàng Dị giới như thể vừa phát hiện ra một sinh vật ngoài hành tinh.

\img{5-1b}

Nàng Sao chổi lên tiếng trước, giọng đầy sững sờ: ``Aki-chan bị làm sao vậy!? Em chưa bao giờ thấy chị ấy cầm sách lâu hơn mười phút! Hôm nay là ngày tận thế à?''

Cậu Tổng - người từ nãy vẫn đang cố gắng làm ngơ - cuối cùng cũng phải ngước lên khỏi màn hình laptop. Anh nhìn đồng hồ, rồi đáp, giọng cũng đầy vẻ bất lực: ``Biết chết liền luôn. Em ấy đã ngồi học được hơn bốn tiếng rồi. Bốn tiếng liên tục, không nghỉ, không uống nước, không đi vệ sinh. Anh bắt đầu nghi ngờ em ấy có phải là Aki như mọi ngày không nữa.''

Nàng Cổ động nhướng mày, tay chống cằm: ``Aki-chan đã biến thành mọt sách rồi à\dots\ Hay là có ai đó đã thay thế cô ấy bằng một bản sao chăm học?''

Câu nói của cô như một nhát dao cứa vào không khí, khiến cả nhóm giật mình. Nhưng còn bất ngờ hơn cả là lời tuyên bố đầy khí phách từ chính nàng Dị giới.

Aki ngẩng đầu lên, chỉnh lại cặp kính mọt sách (từ đâu ra vậy?), giọng đầy hùng hồn: ``Giáo dục sẽ là thứ thống trị thời đại mới! Ai không học sẽ bị bỏ lại phía sau!''

Marine đứng cạnh, mắt mở to, giọng lắp bắp: ``Đến cách nói cũng thay đổi luôn\dots\ Chị ấy bị ai tẩy não vậy? Hay là do đọc sách phát triển bản thân quá nhiều?''

Cậu Tổng gật gù, cảm thấy có gì đó thực sự không ổn: ``\textit{Vãi chưởng\dots} Mặc dù đúng là thế thật, nhưng mà\dots\ thế này là hơi quá rồi đấy. Em ấy đọc sách gì mà thành ra thế nhỉ?''

Nàng Dị giới vẫn tiếp tục cúi đầu vào sách, môi lẩm bẩm không ngừng với những âm thanh lộn xộn, chẳng ai nghe rõ. Thuyền trưởng tiến lại gần, nhón chân, tò mò hỏi: ``Chị lẩm bẩm cái gì vậy? Có công thức nào hay thì chia sẻ với em với.''

Đột nhiên, một câu nói vang lên to và rõ, như thể Aki vừa khám phá ra chân lý của vũ trụ: ``Ba nhân ba bằng\dots\ Chòm Nam Thập Tự!''

\img{5-1c}

Cả căn phòng như chựng lại. Marine đứng gần nhất suýt ngã ngửa, tay bám vào bàn. Kuon và Suisei thì đúng nghĩa là ngã ngửa ra sau, đầu suýt đập vào nhau.

``Không biết chị ấy đang học cái gì, nhưng kiểu này là nước đổ lá khoai thôi!'' Thuyền trưởng vừa đứng dậy vừa nói, giọng đầy hoảng hốt.

Cậu Tổng lồm cồm bò dậy, xoa gáy, lắc đầu ngán ngẩm: ``Thế này thì chữ thầy trả thầy hết rồi còn gì\dots\ Ba nhân ba mà bằng chòm sao thì em ấy nên thi vào khoa thiên văn chứ không phải toán.''

Cả nàng Hải tặc lẫn nàng Sao chổi đều lộ rõ vẻ lo lắng. Suisei đưa tay lên ôm mặt, thở dài: ``Nếu cứ học kiểu này thì có khi em ấy sẽ\dots\ biến thành một nhà bác học điên mất thôi.''

\gif{5-1d}{1}{97}

Trong đầu Suisei và Matsuri bất giác hiện ra một cảnh tượng kỳ quái: Aki đứng trên đỉnh núi, tay cầm một cục đá, mắt sáng rực như vừa tìm ra chân lý cuối cùng của nhân loại.

``Cục đá sau khi đánh bóng thì sáng loáng rồi này!'' nàng Dị giới trong suy nghĩ của hai cô nàng reo lên, giọng đầy phấn khích, ``Không biết là nó cứng cỡ nào nhỉ\dots\ Chắc phải thử mới biết.''

Và rồi, trước sự kinh hoàng của họ, cô thử nghiệm luôn bằng cách\dots\ dùng đầu mình để đập thẳng vào cục đá.

Kết quả: cô nàng gục xuống ngay tại chỗ, trên người hiện ra một chữ ``Chết (死 - Tử)'' màu đỏ chót, như thể vừa bị game over trong một tựa game kinh dị.

Suisei lập tức lắc đầu mạnh, đưa tay xua xua như đang đẩy lùi một con ma: ``Không, không, không! Đừng có tưởng tượng mấy thứ đó nữa! Em mà tưởng tượng nữa chắc em cũng điên mất!''

Marine và Kuon nhìn nhau, cùng chép miệng.

Marine lên tiếng, giọng cố tỏ ra vững vàng, nhưng mắt vẫn còn hơi hoảng: ``Chị ấy sẽ không có đần đến thế đâu!! Đập đá vào đầu thì chẳng khác gì tự sát, Aki-chan có học dốt đến mấy cũng không làm thế được!''

Cậu Tổng gật đầu đồng tình: ``Ừ. Chả ai điên mà lại thử độ cứng của đá bằng cách đập vào đầu mình cả\dots\ Trừ khi em ấy đang đóng phim hành động.''

Cả nhóm đồng loạt thở dài, một tiếng thở dài tập thể, đều như vắt chanh. Nhưng ánh mắt họ vẫn không rời khỏi nàng Dị giới - người vẫn đang chăm chú đọc sách, miệng lẩm nhẩm những công thức toán học pha lẫn tên các chòm sao, như thể chẳng hề biết gì về những gì đang diễn ra xung quanh.

Nhận thấy tình trạng của Aki thực sự bất thường - một cô gái suốt ngày chỉ biết chạy nhảy, cười đùa, bỗng dưng gắn chặt mông vào ghế suốt bốn tiếng đồng hồ - ba cô gái quyết định phải có hành động. Họ đứng thành vòng tròn nhỏ, thì thầm bàn bạc như một băng nhóm đang lên kế hoạch cho một phi vụ lớn.

``Dù sao thì mình cũng nên mau ngăn em ấy lại!'' Suisei lên tiếng, giọng đầy quyết tâm. Đôi mắt sao chổi của cô lóe lên tia hy vọng, như thể đang chuẩn bị cho một ca giải cứu đầy kịch tính.

Matsuri gật đầu lia lịa, tay chống nạnh: ``Đúng đó! Phải đưa cậu ấy về đúng với con người thật của mình! Aki-chan bình thường vui vẻ, hoạt bát, chứ không phải một\dots\ một cái máy học!''

Tuy nhiên, giữa lúc ba người còn đang suy nghĩ đối sách, cậu Tổng lại có một suy nghĩ khác hẳn. Cậu ngồi trên ghế, vắt chéo chân, thản nhiên nói: ``Anh thấy thế nào chứ, học hành chăm chỉ thế cũng tốt mà. Các em cứ lo xa. Để em ấy học đi, biết đâu mai sau đỗ đạt làm quan.''

Mặc kệ lời nói của cậu quản lý, cả ba cô gái liếc nhìn nhau, rồi đồng loạt gật đầu. Một sự đồng thuận tuyệt đối, không cần lời nói. Dù cậu Tổng có nói gì đi nữa, họ vẫn tin rằng phải có cách nào đó kéo Aki khỏi cái bàn học chết tiệt kia!

Họ lập tức bắt tay vào nghĩ ra những trò tiêu khiển nhằm đánh lạc hướng nàng Dị giới, bất kỳ thứ gì đủ sức lay chuyển một tâm hồn đang bị sách giáo khoa bủa vây.

\ct

Suisei là người đầu tiên ra quân. Cô rón rén bước vào bếp, rồi trở ra với một đĩa bánh dâu còn thơm phức, lớp kem trắng mịn, quả dâu đỏ rực nằm trên cùng trông vô cùng hấp dẫn. Đây là món bánh cô vừa làm hôm qua, định bụng dành cho bữa xế chiều cùng mọi người. Nhưng giờ đây, nó trở thành vũ khí tối thượng.

\img{5-1d1}

``Chị vừa làm cái bánh này, Aki-chan, ăn không? Ngon lắm đấy, kem tươi, dâu ngọt, bánh mềm lắm, tất cả những gì tốt nhất cho một buổi chiều học hành mệt mỏi.''

Suisei đặt đĩa bánh xuống cạnh chồng sách, nhìn Aki với đôi mắt đầy hy vọng, như một chiến binh vừa tung ra chiêu cuối.

Nàng Dị giới vẫn cúi gằm, mắt dán vào trang sách, không buồn ngẩng đầu lên. Giọng cô lạnh tanh, đều đều: ``Chị cứ để nó ra kia đi. Em ăn sau, giờ đang đọc đến chương quan trọng.''

Suisei đứng chết trân, tay vẫn cầm đĩa bánh, miệng há hốc. Cô nhìn chiếc bánh dâu hấp dẫn trên tay, rồi nhìn Aki - người mà chỉ một tháng trước còn sẵn sàng bỏ dở mọi việc để lao vào bánh ngọt - và cảm thấy như thể vừa bị một cú đấm thẳng vào lòng tự trọng.

``\textit{Chết thật\dots\ Aki-chan mà còn chối từ bánh ngọt ư? Tận thế thật rồi,}'' cô nghĩ thầm, mặt buồn rười rượi.

\il

Không nản lòng, Matsuri bước lên. Cô nàng lấy từ trong kho ra một quả bóng đá, món đồ chơi ưa thích của Aki trong những buổi chiều nắng đẹp. Cô tâng nhẹ quả bóng lên không trung, đá qua đá lại vài nhịp, rồi lăn nó đến trước mặt Aki, giọng đầy hào hứng:

\img{5-1d2}

``AkiAki! Sao mình không ra ngoài chơi đá bóng đi? Trời đẹp thế này, nắng vàng, gió mát, ở trong phòng suốt phí quá!''

Matsuri vừa nói vừa làm động tác sút bóng, mắt sáng rực. Cô tin rằng với tình yêu bóng đá bất diệt của Aki, món chiêu này chắc chắn sẽ thành công.

Nhưng Aki lần này còn lạnh lùng hơn. Cô ngước mắt lên một giây, liếc nhìn quả bóng, rồi đáp, giọng đầy khinh khỉnh: ``Phí thời gian. Đá bóng có giúp tôi thi đỗ đại học không? Không. Thì thôi.''

Matsuri đứng đơ như một pho tượng, tay vẫn còn giơ lên trong tư thế rủ rê. Quả bóng lăn lăn rồi dừng lại ở góc phòng, lẻ loi, cô độc.

``\textit{Đến cả trò chơi mà Aki từng vui vẻ tham gia cũng bị từ chối sao!? Cô ấy đã thay đổi thật rồi\dots}'' nàng Cổ động nghĩ thầm, lòng tan nát.

\il

Marine cuối cùng cũng quyết định ra tay. Cô bước tới, hai ngón trỏ đưa lên, miệng mở ra định nói: ``A-''

Nhưng còn chưa kịp phát ra âm tiết thứ hai, Aki đã đáp, giọng đều đều như một cái máy: ``Như trên.''

\img{5-1d3}

Nàng Hải tặc đứng sững. Hai ngón tay vẫn giơ lên, môi vẫn mấp máy, nhưng chẳng có âm thanh nào thoát ra. Cô như một diễn viên bị cắt mic giữa sân khấu, với dáng vẻ thảm hại, tội nghiệp, và hơi buồn cười.

``\textit{Ít nhất cũng nghe em nói hết đã chứ\dots\ chị chỉ muốn mời em đi chơi thôi mà\dots}'' Marine nghĩ thầm, mắt bắt đầu rưng rưng.

\img{5-1e}

Cô quay lưng lại, từ từ tiến về phía chiếc tủ đứng cạnh tường. Rồi cô quỳ xuống, áp mặt vào cánh tủ lạnh lẽo, hai tay buông thõng. Cảnh tượng ấy thê thảm đến nỗi Suisei phải chạy lại vỗ vai an ủi: ``Thôi nào, Marine-chan. Không phải lỗi của em. Em ấy bị nhập đồng rồi, không phải do em không hấp dẫn.''

Marine không đáp, chỉ khẽ lắc đầu, mũi vẫn áp vào tủ, giọng nghẹn ngào: ``Em chỉ muốn mời chị ấy đi ăn kem thôi\dots\ thế mà chị ấy còn chưa cho em nói hết câu\dots''

\ct

Cả ba cô gái đứng nhìn Aki, người vẫn đang cật lực học hành, mắt không rời trang sách, tay ghi chép lia lịa. Sự lo lắng trong lòng họ càng lúc càng lớn. Trán họ lấm tấm mồ hôi, không phải vì nóng, mà vì bất lực với cô nàng quyết tâm này.

Matsuri lên tiếng, giọng đầy bất an: ``Không ổn rồi, em ấy cứng đầu quá\dots\ Chúng ta đã thử bánh, thử bóng đá, thử cả\dots\ mời đi chơi, mà vẫn không lay chuyển được. Có khi cô ấy đã bị thay bằng một bản sao chăm học thật rồi.''

Cậu Tổng, vẫn ngồi trên ghế sô-pha với vẻ thư thái đáng ghét, lên tiếng một cách bình thản như đang giảng giải một chân lý hiển nhiên: ``Bình thường thôi. Khi người ta thực sự tập trung học hành hăng say, mấy lời mật ngọt, rủ rê chẳng thấm vào đâu cả. Anh từng trải qua thời kỳ đó rồi.''

Marine lập tức quay phắt lại, mắt nheo lại đầy bực bội: ``Anh không giúp bọn em thì thôi, lại còn đổ thêm dầu vào lửa nữa! Đáng lẽ ra anh phải lo lắng mới đúng! Aki-senpai mà biến thành mọt sách thì văn phòng này mất đi một nửa tiếng cười mất!''

Cậu nhún vai, giọng vẫn tỉnh bơ: ``Thì anh có nói sai đâu. Bọn anh từng dùi mài kinh sử còn kinh khủng hơn thế nhiều. Có người còn thức trắng đêm, tóc rụng nửa đầu, mắt thâm như gấu trúc\dots\ thế này vẫn còn nhẹ chán. Em ấy mới học có bốn tiếng, chưa thấm vào đâu.''

Nàng Hải tặc bực bội giậm chân, mặt đỏ lên. Suisei thì ôm đầu, vẻ mặt khổ sở như vừa ăn phải ớt: ``Em thì không biết thế nào, nhưng nếu cứ như thế này thì Aki-chan sẽ\dots''

Cô không nói hết câu, bởi trong đầu cô và Matsuri, một viễn cảnh khác lại hiện ra đen tối hơn, kinh hoàng hơn, và vô cùng vô lý.

\gif{5-1f}{1}{92}

Trước mắt họ là hình ảnh Aki cầm một quả dừa, khuôn mặt rạng rỡ, đôi mắt sáng lên như thể vừa phát hiện ra một báu vật.

``A, quả dừa này! Hoan hô!'' nàng Dị giới trong tưởng tượng reo lên, giọng đầy phấn khích.

Rồi, không chút do dự, cô dùng chính đầu mình đập thẳng vào quả dừa, với một lực mà cả võ sĩ chuyên nghiệp cũng phải nể nang.

Kết quả? Cô nàng gục xuống ngay tại chỗ, bất động. Trên người lại hiện ra một chữ ``Chết'' màu đỏ chót, y hệt cảnh tượng đập đá lần trước.

Cả hai người - Suisei và Matsuri - đồng loạt rùng mình, suýt ngã khỏi ghế.

Marine đứng nhìn họ với vẻ mặt không thể tin nổi, tay chống nạnh: ``Cách hai người nhìn nhận chị ấy sai vãi cả lề ra ấy! Aki-chan có điên đến mấy cũng không đập dừa vào đầu được! Có khi chỉ nứt đầu thôi chứ chưa chắc chết ngay!''

Cậu Tổng không nhịn được, bật cười. Anh lắc đầu: ``Với cả, làm gì có chuyện đập dừa vào đầu một phát là ngỏm luôn nhỉ? Dừa cũng chỉ là quả dừa, không phải búa tạ. Anh nghĩ là hai đứa nghĩ hơi quá rồi đấy. Có khi Aki-chan chỉ đau thôi.''

Marine quay sang cậu, mắt rực lửa, giọng gần như gào lên: ``\textbf{Vấn đề không phải ở chỗ đó! Em còn thấy sốc hơn khi anh coi chuyện Aki-senpai ôn thi là bình thường đó!}''

Cô bó tay với suy nghĩ khác thường của cậu Tổng quản lý. Thực tế thì cô cũng chẳng biết phải giải thích thế nào, bởi trong thâm tâm, cô hiểu rằng Aki học hành cũng là điều tốt, nhưng cái tốt này lại khiến cô mất đi một người bạn vui vẻ.

Giữa lúc mọi người còn đang bế tắc, Suisei bỗng vỗ tay một cái bốp, mắt sáng lên như vừa tìm ra chân lý cuối cùng của vũ trụ: ``À, chị nghĩ ra cách để đưa em ấy về rồi!''

Matsuri lập tức quay phắt lại, ánh mắt tràn đầy hy vọng, long lanh như sao: ``Cách nào vậy!? Chị định dùng bùa chú hay gọi hồn à?''

Ngược lại, Marine nheo mắt nhìn cô nàng tóc xanh với vẻ đầy ngờ vực. Kinh nghiệm cho thấy, mỗi khi Suisei hứng khởi, thường kết cục chẳng mấy sáng láng.

Nàng Ca sĩ không nói gì. Cô kéo Matsuri lại bàn học ngay bên cạnh Aki, lôi ra hai quyển vở và hai cây bút. Cô đặt chúng xuống, chỉnh tư thế ngồi thẳng lưng, rồi đắc ý tuyên bố: ``Ta cứ học sao cho giỏi hơn em ấy là được! Cạnh tranh là động lực lớn nhất! Nếu thấy người khác học giỏi hơn, tự nhiên Aki-chan sẽ nảy sinh ý chí vươn lên, rồi quên mất việc ``đập đá, đập dừa'' thôi!''

Nàng Cổ động vỗ tay: ``Phải ha! Sao mình không nghĩ ra sớm nhỉ?''

Và rồi, trước sự bàng hoàng của Marine, cả hai cô gái lao ngay vào bàn học như thể bị nhập đồng. Họ cắm mặt vào sách vở, mắt đọc lia lịa, tay viết ghi chép như những cỗ máy được lập trình sẵn.

\gif{5-1g}{1}{12}

Chưa đầy mười giây sau, bàn học bắt đầu rung lên bần bật, nhưng chẳng phải vì động đất, mà vì tốc độ lật sách và viết bài điên cuồng của hai ``chiến binh học tập'' mới toanh.

Marine đứng nhìn, mồ hôi hột lấm tấm trên trán. Cô thì thầm, giọng đầy lo lắng: ``Cách này có tác dụng không đây\dots\ Hay chỉ làm tình hình thêm rối loạn? Giờ có ba người cắm đầu vào sách rồi, lát nữa biến thành bốn người thì sao?''

Cậu Tổng ngồi trên ghế sô-pha, khoanh tay nhìn cảnh tượng với vẻ thích thú. Cậu bình thản đáp, giọng vừa buồn cười vừa hoài niệm: ``Cái này đúng là hơi quá, nhưng mà động cơ của họ cũng tốt mà. Với cả, cảnh này làm anh nhớ đến hồi bọn anh ôn thi vào cấp ba: cũng tranh nhau học, cũng lật sách như lật bánh tráng, cũng viết đến mỏi cả tay.''

Marine giật mình, tò mò hỏi: ``Không phải thi đại học ạ? Em cứ tưởng ở Nhật thi đại học mới căng thẳng cơ.''

Cậu lắc đầu, ánh mắt xa xăm: ``Không. Ở chỗ bọn anh, thi vào cấp ba còn khó hơn thi đại học nhiều. Có những trường tỷ lệ chọi lên đến 1/10, 1/12, thậm chí cao hơn. Nên hồi đó bọn anh học ngày học đêm, ai cũng như ai.''

Marine nhìn cậu với sự đồng cảm sâu sắc, lần đầu tiên cô thấy cậu quản lý không còn vẻ ``lạnh lùng, thờ ơ'' mà có chút gì đó rất người, rất đời. Rồi cô quay sang nhìn ba con người đang cắm cúi vào sách vở: Aki vẫn đều đều lật trang, Suisei và Matsuri thì hùng hục ghi chép như thể đang thi Olympic.

Cô thở dài, giọng trầm xuống: ``Hy vọng là Aki-senpai sẽ trở lại bình thường\dots\ Chứ nếu tình trạng này kéo dài, thì văn phòng chúng ta sẽ biến thành thư viện mất.''

\ct

\begin{center}
  \uline{Một tuần sau.}
\end{center}

Không ai biết chính xác điều gì đã xảy ra trong suốt bảy ngày ấy. Những âm mưu, những chiến thuật ``giải cứu'' Aki khỏi bàn học, những lần Suisei và Matsuri cố gắng ``học giỏi hơn'', tất cả như một mớ hỗn độn không đầu không cuối. Nhưng kết quả cuối cùng\dots\ thì quá rõ ràng, và cũng quá\dots\ kỳ quặc.

Toàn bộ văn phòng \hll bây giờ trông chẳng khác gì một bãi xa-van rộng lớn giữa châu Phi. Những bụi cỏ cao vút, xanh rì, mọc lên từ những kẽ gạch, từ những chiếc ghế, thậm chí từ cả mặt bàn làm việc. Tiếng côn trùng râm ran vọng ra từ khắp các ngóc ngách, tạo nên một bản giao hưởng hoang dã. Gió nhẹ thổi qua, làm những cành cây giả phát ra những tiếng xào xạc đầy bí ẩn, như thể có hàng trăm sinh vật nhỏ bé đang thì thầm điều gì đó.

\img{5-1i1}

Marine và Kuon đứng sững giữa ``khu bảo tồn thiên nhiên'' này, mắt mở to, miệng há hốc. Cậu Tổng đưa tay dụi mắt, rồi dụi lại, nhưng khung cảnh trước mặt vẫn không hề thay đổi.

``Cái quái gì\dots?'' cả hai đồng thanh, giọng đầy bàng hoàng.

Họ nhìn quanh, từng bụi cỏ, từng tán cây nhân tạo, từng đám dây leo quấn quanh quạt trần. Có cảm giác như ai đó đã nhấc bổng cả văn phòng lên và thả xuống giữa thảo nguyên Serengeti.

Và điều bất ngờ nhất, cũng là điều khiến cậu Tổng phải chớp mắt liên tục để chắc chắn rằng mình không bị ảo giác, đó là Suisei và Matsuri. Bằng một cách thần kỳ (hay điên rồ) nào đó, hai cô nàng giờ đây không còn là những cô gái bình thường nữa. Họ đã hóa thành hai con báo đốm.

Những tấm da vàng óng, những đốm đen đặc trưng, đôi mắt sắc như dao, và những bước chân nhẹ nhàng, êm ái, rình rập giữa đám cỏ cao. Trông chẳng khác gì hai mãnh thú đang chuẩn bị vồ mồi.

\img{5-1i}

Cậu Tổng đưa tay dụi mắt lần thứ ba, rồi thở dài: ``\dots\ Thật luôn? Giờ cả văn phòng biến thành thế giới động vật rồi à? Anh có cần đi xin giấy phép nuôi thú hoang không?''

Marine đứng bên cạnh, mắt vẫn trợn tròn, giọng run run: ``NatGeo\footnote{Viết tắt của National Geographic - Kênh truyền hình chuyên về động vật hoang dã và thiên nhiên.} tại gia luôn kìa!? Chỉ còn thiếu ông David Attenborough bước ra thuyết minh thôi!''

Cả hai người chỉ biết đứng nhìn trong im lặng, đầu óc trống rỗng, không biết nên xử lý chuyện này ra sao. Báo đốm Suisei đang nằm phơi nắng trên một tảng đá giả, thỉnh thoảng vẫy đuôi. Báo đốm Matsuri thì đang rón rén bò về phía một bụi cỏ, mắt không rời khỏi\dots\ một chiếc bánh quên để trên bàn từ tuần trước.

Kuon quay sang Marine, giọng đầy bất lực: ``Anh chỉ hỏi thật: có nên gọi sở thú đến không, hay chúng ta cứ để mặc?''

Marine lắc đầu, tay ôm mặt: ``Em không biết nữa\dots\ Nhưng mà hình như Suisei-chan đang làm quen với môi trường sống mới khá tốt thì phải.''

Nhưng biết đâu, chẳng ai dám chắc rằng ``bình thường'' có còn tồn tại trong từ điển của \hll nữa hay không.

%==========================================TRUYỆN PHỤ=========================================
\omk

\pov{Hikawa Kuon}

Ngoài hai con báo đốm (mà tôi chắc kèo chính là Suisei và Matsuri, bởi vì ngoài hai cô nàng đó ra, chẳng ai chịu ngồi học đến nỗi biến hình được), chúng tôi còn nhìn thấy Aki. Cô nàng vẫn đeo nguyên cặp kính mọt sách từ tuần trước, trông như một học giả thực thụ vừa bước ra từ viện nghiên cứu. Cặp kính đã sứt một bên, nhưng cô chẳng hề hay biết.

Và rồi, em ấy hắng giọng, chỉnh kính lại (dù cái kính chẳng còn nằm ngay ngắn trên sống mũi nữa), rồi bắt đầu\dots\ thuyết giảng.

\img{5-1j}

``Báo đốm, còn được biết với tên khoa học là Panthera onca.''

Tôi nheo mắt lại. Ừ thì đúng là tên khoa học thật. Cái này em ấy tra trên mạng cũng ra, không có gì đáng ngạc nhiên. Nhưng mà\dots\ thôi kệ, để xem em ấy tiếp tục nói gì nào.

``Động vật có vú ăn thịt thuộc chi Báo và họ Mèo. Thường được thấy ở Akihabara, môi trường sống tự nhiên của chúng vẫn chưa được xác định.''

\dots\ Khoan đã.

Akihabara? Khu phố điện tử? Nơi mà hàng ngàn otaku tụ tập mua anime figure ấy á? Báo đốm ở Akihabara? Chúng nó lang thang giữa các cửa hàng mô hình gundam và quán mèo café ư?

Tôi quay sang nhìn Marine. Cô nàng đang đứng như trời trồng, mắt mở to, miệng mấp máy nhưng chẳng phát ra âm thanh nào. Rõ ràng, cô cũng đang hoang mang không kém tôi.

Nhưng Aki vẫn chưa dừng lại. Em ấy tiếp tục, giọng đều đều như một giáo sư đang đọc bài giảng sau mười năm nghiên cứu: ``Một số tin rằng chúng có thể dịch chuyển tức thời đến mọi nơi trên Trái Đất, trong khi một số tin rằng đôi chân mới là thứ đã cho chúng có tốc độ vượt trội. Đến giờ thì hai bên vẫn chưa hết tranh cãi về vụ việc này.''

\dots\ Tôi đang nghe cái quái gì thế này?

Dịch chuyển tức thời? Đôi chân siêu tốc? Đây là báo đốm hay là siêu nhân vậy?

Tôi ôm trán, cố gắng tiếp thu cái đống ``kiến thức'' nổ não đó. Không biết em ấy đã học hành kiểu gì, tra cứu ở đâu, nhưng kết quả là một mớ hỗn độn giữa động vật học, khoa học viễn tưởng, và truyện tranh. Đến một đứa trẻ cấp tiểu học còn biết báo đốm sống ở rừng Amazon, chứ không phải Akihabara!

Nhưng tôi vẫn tiếp tục nghe nàng Dị giới giải thích, dù có cảm giác càng nghe thì các tế bào não trong đầu tôi càng lần lượt rời bỏ nhiệm sở. Aki thuyết giảng một cách chậm rãi, trầm bổng, như thể đang truyền đạt một chân lý vĩ đại, nhưng nội dung thì chẳng khác gì đàn gẩy tai trâu, hay đúng hơn, đàn gẩy cho\dots\ báo nghe.

``Loài báo vẫn tiếp tục đưa ra những thắc mắc của chúng về loài người: `\textit{Từ đâu mà có chiến tranh?}' hay `\textit{Đâu là lý do khiến các sĩ tử đi dọn phòng vào buổi tối trước ngày thi?}' và nhiều câu hỏi tu từ như vậy nữa.''

\dots\ Báo đốm có thắc mắc về chiến tranh và chuyện dọn phòng trước ngày thi á!? Chúng là động vật! Chúng quan tâm đến đồ ăn, lãnh thổ, và bạn tình, chứ không phải triết học hay giáo dục!

Tôi định lên tiếng ngắt lời, nhưng Aki đã ngước lên bầu trời (cụ thể là lên trần nhà, nhưng trong đầu em ấy lúc này, có lẽ trần nhà đã biến thành bầu trời xanh với mây trắng), và tiếp tục trình bày, giọng tràn đầy cảm xúc: ``Liệu sẽ có ngày con người có thể giải đáp những câu hỏi đó? Cố lên, con người! Các bạn làm được mà!''

Em ấy giơ nắm đấm lên như một nhà hoạt động xã hội đầy nhiệt huyết.

Rồi em ấy cười nhẹ, lại đẩy cặp kính (lần này suýt đẩy rơi), và kết thúc bài thuyết trình bằng một câu siêu kinh điển: ``P/S: Thực phẩm thiết yếu của chúng là spaghetti bay.''

\dots Vâng.

Con người chúng tôi có thể cho các bạn báo đốm một câu trả lời ngay bây giờ, nhưng sau khi ``thu nạp'' thêm thứ kiến thức nổ não này, chúng tôi cũng không biết phải nói gì nữa.

Spaghetti bay? Mì Ý bay trên trời? Báo đốm ăn mì Ý? Hóa ra chúng là động vật ăn tạp kiểu\dots\ Âu hóa?

Senchou đứng đờ đẫn một lúc, rồi kéo nhẹ áo tôi, thì thầm với giọng đầy lo lắng: ``\hl{Có vẻ như Aki-senpai bị hỏng não rồi. Em nghĩ cần gọi bác sĩ tâm thần, không phải bác sĩ nội khoa.}''

Tôi thở dài ngao ngán, đưa tay xoa trán: ``\hl{Ừ. Học nhiều quá nên hóa điên luôn rồi. Anh đã cảnh báo rồi: cái gì nhiều quá cũng không tốt. Học cũng vậy.}''

Cuối cùng, không còn cách nào khác, tôi đành phải gọi điện cho một người bạn đang làm việc ở bệnh viện tâm thần gần nhất, nhờ họ chuẩn bị một phòng riêng cho Aki; ưu tiên có ô cửa kính quan sát nhìn ra vườn, để em ấy có thể ngắm cây cối và tiếp tục ``nghiên cứu'' về báo đốm.

Còn hai con báo đốm - à không, Suisei và Matsuri mới phải - thì bị tôi lôi về nhà bằng cách\dots\ cột dây thừng vào cổ (họ bảo ``tụi em chạy nhanh lắm, không cột sao được'') để tìm cách thuần hóa, đưa họ trở lại làm người bình thường. Trên đường về, báo đốm Suisei thỉnh thoảng kêu ``meo meo'' nhỏ nhẹ, trong khi báo đốm Matsuri thì nằm lăn ra vỉa hè, đòi phơi bụng.

Quả đúng là ôn thi đại học vất vả thật. Và cái gì nhiều quá cũng không bao giờ là tốt cả, kể cả học, kể cả ăn spaghetti bay, và kể cả việc biến thành động vật hoang dã trong văn phòng. Ừ, và kể cả việc biến hai người đồng nghiệp thành hai con báo đốm vô tri nữa.

\begin{center}{\abv\Large Bonus}\end{center}

Về hai câu hỏi mà loài báo đốm (hoặc là do Aki tự tưởng tượng như vậy) đặt ra, ChatGPT có trả lời như sau:

\img[15cm]{5-1k1}

\img[15cm]{5-1k2}

\end{document}